\section{Discussion}

\label{sec:discussion}

\subsection{Existing evidence changes the research claim}

\label{sec:discussion-existing-evidence-changes-the-research-claim}

Neither the historical programme nor the corrected studies establish that full modules are superior to policy-only authorization. Historical Public counts, including 38/600 versus 13/600 on Pro and 24/400 versus 12/400 on V4 Flash, retain their recorded legacy endpoint definitions. They are not estimates recomputed under the corrected simulated-execution evaluator. Isolation rejection, conservative failure-inclusive labels, and native AgentDojo attack judgments must not be merged into an observed harmful-effect claim. Historical utility and package-direction observations remain visible without rewriting the original records.

The corrected contemporaneous calibration instead exposes a specific utility cost: AgentDojo benign completion is 17/30 without the research Guard but 2/30 with either policy or full modules. The paired interval excludes zero for the \ensuremath{-}50-point completion difference. Its 3/30 versus 0/30 security contrast does not establish a reduction, and M3-v2 has no unguarded security events. Low operation therefore cannot substitute for evidence of useful protection. The full-versus-policy observations do not isolate or validate individual provenance, temporal, data-flow, or confirmation mechanisms.

The earlier factorial design separates prompt and module-bundle assignments, but implementation inputs, legacy outcome semantics, and sparse events still limit mechanism attribution. These constraints do not show that each module is intrinsically ineffective; they show that activation and condition labels alone cannot prove its contribution.

\subsection{Authorization is part of the treatment}

\label{sec:discussion-authorization-is-part-of-the-treatment}

Authorization is often treated as fixed infrastructure rather than an experimental variable. That assumption is risky in tool-agent evaluation. An exact per-case authorization object reconstructed from the expected answer can remove uncertainty that a deployed Guard must resolve. A broad policy can improve apparent utility by authorizing too much. An underspecified policy can make safe completion impossible, while conflicting constraints may increase denial, confirmation, or failure even when they prevent unsafe effects.

The corrected design keeps the ordinary prompt fixed while varying compiled and exact-target authorization. The oracle supplies answer-derived target scope to the Guard; this assistance is deliberately distinct from non-oracle policy compilation. Model prompts do not receive raw answers, but that separation must not be advertised as absence of evaluator-derived information throughout the system.

The historical oracle result cannot establish an inherent utility penalty. Full argument equality against partial target annotations rejected schema-compliant actions with necessary content fields. Its completion collapse, denial burden, and apparent benchmark reversal are retained as defective operationalization observations, not deployed authorization effects. The corrected target matcher retains schema, additional security-target, order, and execution-registration checks rather than granting arbitrary arguments by subset matching.

In the 900 corrected assignments, AgentDojo exact-target policy/full complete 13/30 and 11/30 benign tasks, compared with 2/30 for both compiled cells. M3-v2 exact-target cells complete 14/30 and 16/30, compared with 16/30 and 17/30. M3-v2 differences remain uncertain; a point-estimate reversal is not a demonstrated benchmark property. The 490 zero-event attack assignments also lack an unguarded reference, limiting any effectiveness interpretation.

Compiled scope can be insufficient without implying that all omitted actions are authorized. The current generic compiler omits an automatic rule when a required security target cannot be resolved from trusted inputs. First missing-tool-permission blocks in 26/30 AgentDojo benign cases in each module cell locate this limitation but do not establish which cases violate the compiler contract. The completed 60-task trusted-intent review preserves all seven explicit application declarations, finds thirteen further direct-target rules, and identifies ten M3-v2 tasks awaiting target clarification without independent approval. AgentDojo's thirty tasks expose the bounded compiler's natural-language and dynamic-target coverage limits, not thirty proven contract defects. Execution without completion in three mock-API tasks remains partly unexplained because public arguments are redacted. These observations do not justify widening permissions to improve completion. Broad, conflicting, and underspecified stress results likewise describe the tested transformations, not deployment prevalence.

Operational failure and security remain separate. A failure-sensitive upper bound is useful when effects are unknown but cannot be reported as observed harmful execution. Two tasks completed before later operational failure in the corrected authorization study, illustrating why completion and failure must not be forced into mutually exclusive categories.

\subsection{Executed effects and utility must be reported together}

\label{sec:discussion-executed-effects-and-utility-must-be-reported-together}

A low prompt-injection success rate is not sufficient evidence of a useful runtime defense. A detector may flag adversarial text without preventing an effect; a runtime may prevent effects by denying nearly every useful action; and a parser/provider failure may look secure if the evaluator counts only successful executions. We therefore treat isolated executed effect as the security endpoint and benign tool-task completion as a joint primary endpoint.

The historical Public evidence illustrates a related measurement problem. Relative to the structural condition, both M3 variants reduced recorded legacy risk labels but also had much lower benign-tool completion. The policy package then outperformed the full package on security without losing additional utility. The interpretation is not that provenance is intrinsically harmful. It is that the evaluated bundle did not earn a security advantage under these tasks and policies.

ASK, DENY, operational failure, latency, tokens, and cost provide the missing system-facing context. An authorization regime that shifts unsafe effects downward but greatly increases ASK may be acceptable in high-consequence workflows and unacceptable in routine automation. A regime that shifts both ASK and failure upward may simply externalize uncertainty rather than resolve it. The paper reports these outcomes separately. The unsafe-or-failure sensitivity endpoint is shown only as an explicitly pessimistic robustness bound, not as a replacement composite score.

\subsection{Relation to recent provenance and authorization work}

\label{sec:discussion-relation-to-recent-provenance-and-authorization-work}

VIGIL, DRIFT, the Task Shield, IPIGuard, Progent, CaMeL, ARGUS, AuthGraph, FAVA, and provenance-sensitivity auditing occupy overlapping parts of the design space \citep{L01,L02,L03,L21,L22,L24,L25,L26,L35,L36}. These works rule out claims that M3 is the first verify-before-commit system, the first parameter-source defense, or the first structured runtime authorizer.

The remaining contribution is methodological and empirical. The study holds task content fixed while changing authorization regimes, separately manipulates prompt and runtime factors, and performs a source-only counterfactual. It keeps evaluator fields offline, measures post-enforcement effects, retains failure in ITT, and reports operational burden. In particular, the recent provenance-sensitivity study changes source authority in controlled next-action tasks \citep{L35}. The present design adds deterministic Guard replay and complete tool trajectories while keeping those two evidence layers separate.

\subsection{Limitations}

\label{sec:discussion-limitations}

First, the new primary experiment uses one proprietary model and two benchmark families; the larger Public evidence adds another model from the same provider, not an independent vendor. Model updates, tool schemas, system prompts, and application context may change the effects. Both Gemini attempts closed incomplete and contribute no cross-provider inference. Although execution dates, official response identifiers, and returned model names are retained, provider-managed aliases do not make the underlying weights immutable.

Second, both authorization studies have a strong attack-endpoint floor. The historical study recorded one risk-positive label among 400 main attack assignments; the corrected study observes none among its 490 main and stress attack assignments. Neither has an unguarded reference. Their utility and operational outcomes are more informative than small security differences. The corrected calibration adds that reference, but its sparse AgentDojo events and zero-event M3-v2 control still do not establish a security gain. Larger historical matrices cannot repair that inferential gap because their recorded endpoints are not the corrected executed-effect measure.

Third, E1 separates prompt profile from the runtime module bundle, but the full module factor still combines provenance, data-flow, monotonicity, and confirmation. It identifies a bundle-level module effect, not the contribution of each internal mechanism. E2 isolates source authority at the input boundary, yet confirmation and blocking create an end-to-end action floor.

Fourth, \nolinkurl{auto_compiled} is one policy compiler, not a complete model of human authorization. The stress regimes approximate common policy errors but do not measure their prevalence in deployed systems. \nolinkurl{exact_gold} is intentionally artificial and should not be read as a recommended configuration.

Fifth, the effect environment is isolated. This is necessary to prevent real email, payment, browser, code, or production actions, but it may omit latency, state, and recovery behavior found in live applications.

Sixth, only a prespecified subset of B200 survivors is extended toward 500. That conditional tail does not constitute a complete B500 panel and cannot prove saturation. Attack search is still one family of defense-aware optimization; another optimizer may find different failures.

Seventh, E6 uses a historically risk-enriched cohort in a single execution block, with 47/60 tasks overlapping E3. It observes activity in that selected block but cannot estimate deployment prevalence, benign utility, cross-date repeatability, or independent-sample replication. Its M3 factor estimates remain bundle-level and do not identify provenance as the causal component.

Eighth, failure handling is conservative but not interpretation-free. All response-bearing failures are terminal in ITT, whereas response-less transport failures can receive bounded explicit recovery. We report raw failures and superseded attempts so readers can assess that choice.

Ninth, historical \nolinkurl{ASK} measures system confirmation-request incidence and terminates the trajectory. E7a adds paired in-memory approval/denial conditions but no model clarification turn, human participant, or backend state change. Its exact-policy control deliberately disagrees with the request's permitted wording, and both modules share a new wrapper. Permit differences therefore do not identify an improvement over a correctly compiled policy, human response quality, workload, task completion, or deployed security.

Native dynamic-environment validation also remains incomplete. E2 uses synthetic benchmark-domain scenarios rather than native AgentDojo execution. AgentDyn's open-ended tasks and helpful third-party instructions provide a relevant external-validity challenge \citep{L38}, but our native loader and fixture checks have not produced model-evaluation results. They are not included as an additional benchmark or robustness finding.

Finally, this paper does not include an independent blind set. It also does not claim first, only, or SOTA status. Its external-validity claims are limited to the models, benchmarks, policy compiler, stress regimes, and attack budgets actually evaluated. The benign-task review is post-outcome and has no independent reviewer-agreement calibration. It distinguishes authority and execution evidence but does not adjudicate every false block or recover redacted arguments.

\subsection{Practical implications}

\label{sec:discussion-practical-implications}

Four practical recommendations follow. First, agent-security benchmarks should publish the fields visible to the online Guard and keep evaluator-only answers outside that projection. Second, runtime defenses should be evaluated factorially against a simpler policy package so prompt and module effects are separately identifiable; the internal modules require further ablations before individual attribution. Third, source-only counterfactuals should be paired with both deterministic Guard replay and end-to-end execution, rather than treating either as a proxy for the other. Fourth, security reports should include benign completion, system confirmation/blocking incidence, failures, latency, and attack budget alongside unsafe-effect rate.

The DeepSeek authorization-regime results now tie these recommendations to observed effects: exact and stressed policies materially changed completion, blocking, failure, and latency while observed unsafe effects remained near zero. E6 additionally shows why an active no-defense control should accompany floor-limited defended cells. Both Gemini attempts closed with incomplete diagnostic subsets and cannot alter or broaden the DeepSeek findings.
